% GENERATED FILE. Edit canonical lessons, then run npm run generate:books.
% canonical-source-hash: fnv1a64:704aadadd6096dfd
% canonical-lessons: TA-C492-tantai, TA-C492-tay, TA-C492-cakotaran, TA-C492-cakotari, TA-C492-kataikkutti

\chapter{Mother and Father}
\label{ch:ta-mother-and-father}
\hlchaptermodality{492}

\begin{chapteropening}
\textbf{By the end of this chapter:} \emph{I can use the written words for mother, father, brother and sister.}

Complete the last lesson of chapter 492: I can use the written words for mother, father, brother and sister.
\end{chapteropening}

\section[\texorpdfstring{\ta{தந்தை} (tantai)}{தந்தை (tantai)}]{\ta{தந்தை} (tantai) — one’s father; a father}

\label{lesson:TA-C492-tantai}



\begin{quote}
Before the new one: say the Tamil for to groan, then the Tamil for surgery.
\end{quote}

\subsection*{You'll want to know: \ta{தந்தை}}
\textbf{\ta{தந்தை}} — \emph{tantai} — \textquotedblleft{}one’s father; a father\textquotedblright{}.

\begin{quote}
\ta{என்} \ta{தந்தை} \ta{ஒரு} \ta{ஆசிரியர்}.
\end{quote}

\begin{quote}
\emph{eṉ tantai oru āciriyar.}
\end{quote}

\begin{quote}
My father is a teacher.
\end{quote}

Written and formal; at home \ta{அப்பா}.

\begin{cousinweb}
\textbf{\ta{தந்தை}} (\emph{tantai}) is native Dravidian (DEDR 3067); the same entry has \ta{எந்தை} \emph{entai}, \textquotedblleft{}my father\textquotedblright{}, in old Tamil.
\end{cousinweb}

\begin{grammarlens}[title={What you've built}]
One more word for describing what you lived and felt.
\end{grammarlens}

\subsection*{Your turn}
\begin{itemize}
\raggedright
  \item \emph{Say it:} \emph{tantai}
  \item \emph{Say it:} \emph{tantai}, once more
  \item \emph{Say it:} \emph{eṉ tantai oru āciriyar.} — My father is a teacher.
  \item \emph{Recall:} say the Tamil for to groan, then the Tamil for surgery, then say \emph{tantai} again
\end{itemize}

\begin{tcolorbox}[breakable,colback=teal!4,colframe=teal!35!black,title={Before you move on}]
What does \ta{தந்தை} mean? (\textquotedblleft{}One’s father; a father\textquotedblright{}.) Say it once more.
\end{tcolorbox}
\section[\texorpdfstring{\ta{தாய்} (tāy)}{தாய் (tāy)}]{\ta{தாய்} (tāy) — a mother}

\label{lesson:TA-C492-tay}



\begin{quote}
Before the new one: say the Tamil for surgery, then the Tamil for one’s father.
\end{quote}

\subsection*{You'll want to know: \ta{தாய்}}
\textbf{\ta{தாய்}} — \emph{tāy} — \textquotedblleft{}a mother\textquotedblright{}.

\begin{quote}
\ta{தாய்} \ta{தன்} \ta{குழந்தையைத்} \ta{தூக்கினாள்}.
\end{quote}

\begin{quote}
\emph{tāy taṉ kuḻantaiyait tūkkiṉāḷ.}
\end{quote}

\begin{quote}
The mother picked up her child.
\end{quote}

Written and formal; at home \ta{அம்மா}. \ta{தாய்மொழி}, a mother tongue.

\begin{cousinweb}
\textbf{\ta{தாய்}} (\emph{tāy}) is native Dravidian (DEDR 364), in an entry of words for mother such as \ta{ஆய்} \emph{āy}.
\end{cousinweb}

\begin{grammarlens}[title={What you've built}]
One more word for describing what you lived and felt.
\end{grammarlens}

\subsection*{Your turn}
\begin{itemize}
\raggedright
  \item \emph{Say it:} \emph{tāy}
  \item \emph{Say it:} \emph{tāy}, once more
  \item \emph{Say it:} \emph{tāy taṉ kuḻantaiyait tūkkiṉāḷ.} — The mother picked up her child.
  \item \emph{Recall:} say the Tamil for surgery, then the Tamil for one’s father, then say \emph{tāy} again
\end{itemize}

\begin{tcolorbox}[breakable,colback=teal!4,colframe=teal!35!black,title={Before you move on}]
What does \ta{தாய்} mean? (\textquotedblleft{}A mother\textquotedblright{}.) Say it once more.
\end{tcolorbox}
\section[\texorpdfstring{\ta{சகோதரன்} (cakōtaraṉ)}{சகோதரன் (cakōtaraṉ)}]{\ta{சகோதரன்} (cakōtaraṉ) — a brother}

\label{lesson:TA-C492-cakotaran}



\begin{quote}
Before the new one: say the Tamil for one’s father, then the Tamil for a mother.
\end{quote}

\subsection*{You'll want to know: \ta{சகோதரன்}}
\textbf{\ta{சகோதரன்}} — \emph{cakōtaraṉ} — \textquotedblleft{}a brother\textquotedblright{}.

\begin{quote}
\ta{எனக்கு} \ta{இரண்டு} \ta{சகோதரர்கள்} \ta{இருக்கிறார்கள்}.
\end{quote}

\begin{quote}
\emph{eṉakku iraṇṭu cakōtararkaḷ irukkiṟārkaḷ.}
\end{quote}

\begin{quote}
I have two brothers.
\end{quote}

Formal; at home \ta{அண்ணன்} or \ta{தம்பி}. Plural \ta{சகோதரர்கள்}: the -\ta{அன்} becomes -\ta{அர்} before -\ta{கள்}.

\begin{cousinweb}
\textbf{\ta{சகோதரன்}} (\emph{cakōtaraṉ}) is borrowed from Sanskrit \emph{sahōdara}, \textquotedblleft{}born of the same womb\textquotedblright{} (CDIAL 13312), with the Tamil -\ta{அன்} of a man.
\end{cousinweb}

\begin{grammarlens}[title={What you've built}]
One more word for describing what you lived and felt.
\end{grammarlens}

\subsection*{Your turn}
\begin{itemize}
\raggedright
  \item \emph{Say it:} \emph{cakōtaraṉ}
  \item \emph{Say it:} \emph{cakōtaraṉ}, once more
  \item \emph{Say it:} \emph{eṉakku iraṇṭu cakōtararkaḷ irukkiṟārkaḷ.} — I have two brothers.
  \item \emph{Recall:} say the Tamil for one’s father, then the Tamil for a mother, then say \emph{cakōtaraṉ} again
\end{itemize}

\begin{tcolorbox}[breakable,colback=teal!4,colframe=teal!35!black,title={Before you move on}]
What does \ta{சகோதரன்} mean? (\textquotedblleft{}A brother\textquotedblright{}.) Say it once more.
\end{tcolorbox}
\section[\texorpdfstring{\ta{சகோதரி} (cakōtari)}{சகோதரி (cakōtari)}]{\ta{சகோதரி} (cakōtari) — a sister}

\label{lesson:TA-C492-cakotari}



\begin{quote}
Before the new one: say the Tamil for a mother, then the Tamil for a brother.
\end{quote}

\subsection*{You'll want to know: \ta{சகோதரி}}
\textbf{\ta{சகோதரி}} — \emph{cakōtari} — \textquotedblleft{}a sister\textquotedblright{}.

\begin{quote}
\ta{என்} \ta{சகோதரி} \ta{வெளிநாட்டில்} \ta{வேலை} \ta{செய்கிறாள்}.
\end{quote}

\begin{quote}
\emph{eṉ cakōtari veḷināṭṭil vēlai ceykiṟāḷ.}
\end{quote}

\begin{quote}
My sister works abroad.
\end{quote}

Formal; at home \ta{அக்கா} or \ta{தங்கை}.

\begin{cousinweb}
\textbf{\ta{சகோதரி}} (\emph{cakōtari}) is borrowed from Sanskrit \emph{sahōdara} (CDIAL 13312), with the -\ta{இ} of a woman.
\end{cousinweb}

\begin{grammarlens}[title={What you've built}]
One more word for describing what you lived and felt.
\end{grammarlens}

\subsection*{Your turn}
\begin{itemize}
\raggedright
  \item \emph{Say it:} \emph{cakōtari}
  \item \emph{Say it:} \emph{cakōtari}, once more
  \item \emph{Say it:} \emph{eṉ cakōtari veḷināṭṭil vēlai ceykiṟāḷ.} — My sister works abroad.
  \item \emph{Recall:} say the Tamil for a mother, then the Tamil for a brother, then say \emph{cakōtari} again
\end{itemize}

\begin{tcolorbox}[breakable,colback=teal!4,colframe=teal!35!black,title={Before you move on}]
What does \ta{சகோதரி} mean? (\textquotedblleft{}A sister\textquotedblright{}.) Say it once more.
\end{tcolorbox}
\section[\texorpdfstring{\ta{கடைக்குட்டி} (kaṭaikkuṭṭi)}{கடைக்குட்டி (kaṭaikkuṭṭi)}]{\ta{கடைக்குட்டி} (kaṭaikkuṭṭi) — the youngest child}

\label{lesson:TA-C492-kataikkutti}



\begin{quote}
Before the new one: say the Tamil for a brother, then the Tamil for a sister.
\end{quote}

\subsection*{You'll want to know: \ta{கடைக்குட்டி}}
\textbf{\ta{கடைக்குட்டி}} — \emph{kaṭaikkuṭṭi} — \textquotedblleft{}the youngest child\textquotedblright{}.

\begin{quote}
\ta{வீட்டில்} \ta{கடைக்குட்டி} \ta{என்பதால்} \ta{அவளுக்குச்} \ta{செல்லம்} \ta{அதிகம்}.
\end{quote}

\begin{quote}
\emph{vīṭṭil kaṭaikkuṭṭi eṉpatāl avaḷukkuc cellam atikam.}
\end{quote}

\begin{quote}
As the youngest in the family she is much pampered.
\end{quote}

\ta{கடை} \textquotedblleft{}the end\textquotedblright{} + \ta{குட்டி} \textquotedblleft{}a little one\textquotedblright{}.

\begin{cousinweb}
\textbf{\ta{கடைக்குட்டி}} (\emph{kaṭaikkuṭṭi}) is a native compound of \ta{கடை}, \textquotedblleft{}the end, the last\textquotedblright{}, and \ta{குட்டி} \emph{kuṭṭi}, \textquotedblleft{}a young one\textquotedblright{} (DEDR 1670).
\end{cousinweb}

\begin{grammarlens}[title={What you've built}]
One more word for describing what you lived and felt.
\end{grammarlens}

\subsection*{Your turn}
\begin{itemize}
\raggedright
  \item \emph{Say it:} \emph{kaṭaikkuṭṭi}
  \item \emph{Say it:} \emph{kaṭaikkuṭṭi}, once more
  \item \emph{Say it:} \emph{vīṭṭil kaṭaikkuṭṭi eṉpatāl avaḷukkuc cellam atikam.} — As the youngest in the family she is much pampered.
  \item \emph{Recall:} say the Tamil for a brother, then the Tamil for a sister, then say \emph{kaṭaikkuṭṭi} again
\end{itemize}

\begin{tcolorbox}[breakable,colback=teal!4,colframe=teal!35!black,title={Before you move on}]
What does \ta{கடைக்குட்டி} mean? (\textquotedblleft{}The youngest child\textquotedblright{}.) Say it once more.
\end{tcolorbox}
